\documentclass[letterpaper,journal]{IEEEtran}

% Core math and IEEE-compatible utilities
\usepackage{amsmath,amsfonts}
\usepackage{array}
\usepackage[caption=false,font=normalsize,labelfont=sf,textfont=sf]{subfig}
\usepackage{textcomp}
\usepackage{stfloats}
\usepackage{url}
\usepackage{verbatim}
\usepackage{graphicx}
\usepackage{cite}

% Packages required by the migrated ACM sample content
\usepackage{multirow}
\usepackage{colortbl}
\usepackage{booktabs}
\usepackage{tikz}
\usetikzlibrary{shapes.geometric, arrows.meta, positioning, fit, backgrounds, calc}

\hyphenation{op-tical net-works semi-conduc-tor IEEE-Xplore}

% -----------------------------------------------------------------------------
% Revised 2026-06-26: Main narrative rewritten to HLC-CQR.
% Old Gaussian-B / residual-diffusion / incident-regime content preserved as
% diagnostic/background sections.  All HLC-CQR numbers sourced from 8-seed
% STID experiments on PeMS08/PeMS04.  No new bib keys introduced.
% -----------------------------------------------------------------------------

\begin{document}

\title{HLC-CQR: Hierarchical Conformal Calibration for Reliable Traffic Forecasting}

\author{Ziyu~Wang, Kaijie~Xu, and Lin~Qiu%
\thanks{Ziyu Wang, Kaijie Xu, and Lin Qiu are with Zhejiang University, Hangzhou, Zhejiang, China.}%
\thanks{E-mail: yan.24@intl.zju.edu.cn; Kaijie.23@intl.zju.edu.cn; linqiu@intl.zju.edu.cn.}%
\thanks{The authors contributed equally to this work.}}

% The paper headers. Update the journal/volume fields before submission.
\markboth{Journal Manuscript Draft}%
{Wang \MakeLowercase{\textit{et al.}}: HLC-CQR: Hierarchical Conformal Calibration for Traffic Forecasting}

\maketitle

\begin{abstract}
Reliable probabilistic traffic forecasting requires calibrated prediction
intervals that are simultaneously sharp, conditionally valid, and robust to
scarce calibration data.  Modern deterministic backbones such as STID
already supply strong frozen point forecasts; the remaining challenge is
post-hoc interval calibration that does not disturb the frozen mean.  Split
conformal prediction provides marginal coverage but yields coarse,
overly wide intervals.  Load-conditioned and time-of-day-conditioned
conformalized quantile regression (CQR-load, Mondrian conformal) sharpens
intervals by conditioning on traffic regimes, but fine-grained groups such
as node$\times$horizon$\times$load$\times$TOD are chronically under-sampled on
standard calibration splits, causing erratic quantile estimates and
coverage collapse.

We propose \emph{HLC-CQR} (Hierarchical Load-Conditioned Conformalized
Quantile Regression), a traffic-structured hierarchical conformal calibration
method that resolves the granularity--sample-efficiency trade-off.  HLC-CQR
organizes traffic-predictive features into a five-level group hierarchy
(G0\,global $\to$ G1\,horizon $\to$ G2\,node$\times$horizon $\to$
G3\,node$\times$horizon$\times$load $\to$
G4\,node$\times$horizon$\times$load$\times$TOD) and applies
sample-size-adaptive shrinkage
$Q^{*}_{g}=\lambda_{g}Q_{g}+(1-\lambda_{g})Q^{*}_{\text{parent}(g)}$
with $\lambda_{g}=n_{g}/(n_{g}+\tau)$ to borrow strength from parent
groups when a fine-level cell is sparse.  The hyperparameter $\tau=20$ is
pre-registered; $\tau=0$ recovers naive Mondrian conformal and $\tau\to\infty$
collapses to the coarse parent.  HLC-CQR is post-hoc and never modifies the
frozen mean.  It makes no new conformal coverage theorem; CQR-load is a
traffic-specific Mondrian/group-conditional conformal instantiation, and
HLC-CQR's novelty lies in the traffic hierarchy and shrinkage mechanism.

On PeMS08 and PeMS04 with 8 seeds, HLC-CQR achieves Winkler scores of
$87.69$ and $111.22$, outperforming standard conformal and CQR-load baselines,
and beats CSDI by $\Delta{=}{-9.38}$ (PeMS08) and $\Delta{=}{-9.18}$
(PeMS04) on the Winkler interval score.  Under 5\% calibration data, HLC-CQR
degrades gracefully (Winkler~$91.70$) while naive CQR-load$\times$TOD
deteriorates to~$95.33$ on PeMS08.  The method is portable across STID,
STAEformer, AGCRN, and PDFormer backbones, and is approximately $250\times$
faster than localized conformal at inference.  We do not claim CRPS
superiority over generative baselines such as CSDI; our honest scope is
interval sharpness and calibration robustness.
\end{abstract}

\begin{IEEEkeywords}
Probabilistic traffic forecasting, conformal prediction, conformalized quantile regression, hierarchical calibration, interval sharpness, intelligent transportation systems.
\end{IEEEkeywords}

\section{Introduction}

Accurate and calibrated uncertainty estimates are essential in intelligent
transportation systems (ITS), where operators must plan not only around the
most likely future traffic state but also around deviations caused by
incidents, weather, and demand surges.  State-of-the-art deterministic
backbones---STGCN, DCRNN, GWNet, AGCRN, PDFormer, STID, and
STAEformer~\cite{yu2018stgcn,li2018dcrnn,wu2019graphwavenet,bai2020agcrn,jiang2023pdformer,shao2022stid,liu2023staeformer}---have
substantially reduced point-prediction error.  The complementary problem is
to provide \emph{calibrated prediction intervals} around those frozen
deterministic forecasts with minimal additional computation and without
perturbing the point estimate.

\textbf{Split conformal prediction} offers a post-hoc route to marginal
coverage guarantees: apply a holdout calibration set to find the quantile of
residual nonconformity scores and center an interval around the frozen mean.
This is architecture-agnostic and computationally inexpensive, but it
produces a \emph{global} quantile that ignores traffic regime variation.  In
practice, the resulting interval is wide everywhere because it must cover
heavy-tailed drop events that are a small fraction of all windows.

\textbf{CQR-load and Mondrian conformal} sharpen intervals by conditioning
on traffic regime features (load level, time of day).  Load-conditioned
conformalized quantile regression (CQR-load) partitions the calibration set
by predicted load bin and computes a separate conformal quantile per
partition.  This substantially reduces interval width for the dominant
moderate-load regime.  However, fine-grained partitions---by
node$\times$horizon$\times$load$\times$TOD---become severely under-sampled
on standard 20\% calibration splits.  Sparse partitions yield high-variance
quantile estimates and can produce out-of-distribution coverage behavior,
negating the sharpness benefit.

\textbf{This paper introduces HLC-CQR}, a traffic-structured hierarchical
conformal calibration framework.  HLC-CQR defines a five-level group
hierarchy anchored in the natural structure of traffic forecasting (sensors,
horizons, load levels, time of day) and applies sample-size-adaptive
shrinkage across the hierarchy so that sparse fine-level cells borrow
calibration strength from their ancestors.  The result is an interval
calibration method that is simultaneously sharper than global conformal,
more stable than naive fine-grained Mondrian, and essentially free at
inference (O(1) lookup rather than kNN search).

\textbf{Relation to prior work.}
Closely related generic residual-calibration and conformal regression work
includes CARD, RDIT, FALDA/DMRR, and diffusion residual
modeling~\cite{han2022card,lai2025rdit,wang2025falda,qi2026diffusionresidual}.
These establish that a strong deterministic predictor plus a post-hoc residual
layer can reach competitive probabilistic scores.  Our work does not claim a
new conformal coverage theorem; CQR-load is already a traffic-specific
Mondrian/group-conditional conformal instantiation.  The novelty of HLC-CQR
lies in the traffic-semantic hierarchy and the shrinkage mechanism that
addresses the granularity--sample-efficiency trade-off inherent in fine-grained
conditioning.  We further provide a diagnostic comparison with diffusion
residuals (CSDI~\cite{tashiro2021csdi}, PriSTI~\cite{liu2023pristi},
DiffSTG~\cite{wen2023diffstg}) as background and limitation context, not
as a competing primary framing.

\textbf{Contributions:}
\begin{itemize}
    \item \textbf{Frozen-backbone residual calibration benchmark.}  We
    establish a controlled, same-backbone evaluation protocol for post-hoc
    interval calibration where the frozen deterministic mean is preserved by
    construction.  All residual calibration methods share identical frozen
    point forecasts, scalers, and splits, enabling fair isolation of
    calibration quality from point-forecast improvements.
    \item \textbf{HLC-CQR algorithm.}  We propose hierarchical
    load-conditioned CQR with a five-level traffic hierarchy
    (G0--G4) and sample-size-adaptive shrinkage
    $Q^{*}_{g}=\lambda_{g}Q_{g}+(1-\lambda_{g})Q^{*}_{\text{parent}(g)}$,
    $\lambda_{g}=n_{g}/(n_{g}+\tau)$, $\tau=20$ pre-registered.
    This resolves the granularity--sample-efficiency trade-off that
    naive Mondrian conformal fails to address.
    \item \textbf{Mechanism validation.}  We validate the shrinkage
    mechanism through $\tau$ ablations across a grid
    $\{0, 5, 10, 20, 50, 100, 200, \infty\}$, hierarchy ablations
    across G0--G4, and a sample-efficiency stress test at calibration
    ratios $\{100\%, 50\%, 25\%, 10\%, 5\%\}$, confirming that HLC-CQR
    degrades gracefully while naive fine-grained Mondrian collapses.
    \item \textbf{Empirical validation.}  On PeMS08/PeMS04 with 8 seeds,
    multi-$\alpha$ ($\{0.05, 0.10, 0.20\}$), and external CSDI Winkler
    anchors, HLC-CQR achieves best Winkler scores of $87.69$/\,$111.22$
    and beats CSDI by $\Delta{=}{-9.38}$/\,$-9.18$ on the interval
    score.  We do not claim CRPS superiority over CSDI.
    \item \textbf{Portability and defensive evidence.}  HLC-CQR is
    ported to STID, STAEformer, AGCRN, and PDFormer backbones; compared
    against localized conformal (HLC wins PeMS08; localized slightly wins
    PeMS04 but at ${\approx}250\times$ higher inference cost); and
    evaluated with a coverage-safety variant (HLC-safe, PICP
    $0.898/0.896$) with negligible Winkler cost.  Gaussian-B and
    residual diffusion are retained as diagnostic baselines for background
    context, not as core contributions.
\end{itemize}

\section{Related Work}

\subsection{Deterministic Spatio-Temporal Traffic Forecasting}
Traffic forecasting has evolved from statistical models toward deep
spatio-temporal learning over urban sensor networks.  Graph-based models
such as STGCN, DCRNN, GWNet, and AGCRN encode road-network
dependencies~\cite{yu2018stgcn,li2018dcrnn,wu2019graphwavenet,bai2020agcrn},
while attention and Transformer-based architectures such as GMAN, PDFormer,
and STAEformer further improve long-range temporal and spatial
modeling~\cite{zheng2020gman,jiang2023pdformer,liu2023staeformer}.  STID
provides a strong MLP-based graph-free baseline~\cite{shao2022stid}.  In
our work, these serve as frozen mean predictors.  A key requirement is that
these backbones must be reproduced at literature-level accuracy; otherwise,
calibration quality is confounded by a weak mean.

\subsection{Probabilistic Forecasting and Diffusion Models}
Diffusion-based models such as CSDI, TimeGrad, PriSTI, DiffSTG, and
SpecSTG extend point prediction to generative predictive
distributions~\cite{tashiro2021csdi,rasul2021timegrad,liu2023pristi,wen2023diffstg,lin2024specstg}.
A closely related line of work including CARD, RDIT, FALDA/DMRR, and
diffusion residual modeling for long-term
forecasting~\cite{han2022card,lai2025rdit,wang2025falda,qi2026diffusionresidual}
decouples a strong deterministic point estimator from residual uncertainty
modeling, establishing that calibrated residual distributions can substantially
improve probabilistic scores without disturbing the deterministic mean.  We
treat this paradigm as an experimental scaffold and diagnostic reference
rather than a methodological claim, because our main contribution is the
HLC-CQR conformal calibration layer.

\subsection{Conformal Prediction and Interval Calibration}
Post-hoc conformal calibration converts any point predictor into a
valid marginal-coverage interval estimator using a held-out calibration set.
Mondrian/group-conditional conformal extends this to conditional coverage
by partitioning the calibration set by group, at the cost of larger
per-group quantile estimation variance.  Conformalized quantile regression
(CQR) uses estimated quantile functions as the base predictor, yielding
sharper intervals than residual conformal.  Localized conformal prediction
further conditions on local neighborhoods in feature space.  Our HLC-CQR
builds on CQR-load (a traffic-specific Mondrian CQR instantiation) and adds
hierarchical shrinkage to manage sparse fine-grained calibration groups.

\subsection{External Semantics for Traffic Forecasting}
Language models and semantic urban data can encode weather, incidents,
calendar events, and regional context.  Earlier versions of this work
treated semantic prompting as a central mechanism.  Our current evidence
does not support a universal semantic-improvement claim; no-semantic
variants are competitive or stronger on aggregate metrics.  Semantics are
therefore restricted to an optional conditioning source in diagnostic
analysis.

\section{Method}

\subsection{Problem Formulation and Notation}
Let $X_{t-L+1:t}\in\mathbb{R}^{L\times N\times F}$ denote historical
traffic observations over $N$ sensors, $L$ input steps, and $F$ features.
The goal is to predict $Y_{t+1:t+H}\in\mathbb{R}^{H\times N\times 1}$ and
a calibrated prediction interval $[\hat{L},\hat{U}]$ at nominal coverage
level $1-\alpha$ (default $\alpha=0.10$).  We evaluate point metrics
(MAE/RMSE) and interval metrics (PICP at level $1-\alpha$, MPIW,
Winkler$_{90}$, and CRPS) after inverse transformation to the physical
traffic scale.

\subsection{Frozen Mean Predictor}
A deterministic backbone $f_\phi$ produces a point forecast
\begin{equation}
    \hat{\mu} = f_\phi(X_{t-L+1:t}),
\end{equation}
where $\phi$ is frozen throughout calibration.  The backbone may be STID,
AGCRN, STAEformer, PDFormer, or another validated deterministic
forecaster.  All calibration methods operate on the residual
$R = Y - \hat{\mu}$ in the normalized space and produce intervals centered
on the backbone prediction: the frozen mean is never modified.

\subsection{Split Conformal and CQR-Load Baselines}
\label{sec:cqr_baseline}

\textbf{Global split conformal.}  Given a calibration set
$\mathcal{D}_\text{cal}=\{(x_i,y_i)\}_{i=1}^{n}$, compute conformity
scores $s_i=|y_i-\hat{\mu}(x_i)|$ and set the interval half-width to the
$(1-\alpha)(1+1/n)$ empirical quantile of $\{s_i\}$.  This provides a
marginal coverage guarantee but produces a single global quantile that does
not adapt to traffic regime.

\textbf{CQR-load.}  Partition the calibration set by a traffic regime
feature (e.g., discretized predicted load) into groups $\{G_k\}$.  Within
each group $G_k$, compute a local conformal quantile $Q_k$ from
$\{s_i : i\in G_k\}$.  At test time, assign each sample to its group and
apply $Q_k$.  This is a direct instantiation of Mondrian/group-conditional
conformal prediction and inherits marginal coverage within each group.
CQR-load$\times$TOD further conditions on time-of-day, yielding finer but
sparser groups.

\textbf{The granularity--sample-efficiency trade-off.}  Finer groups yield
sharper adapted intervals in the large-data regime, but sparse groups
produce high-variance quantile estimates that can under- or over-cover.  In
a standard 20\% calibration split for PeMS08 ($N{=}170$ sensors,
$H{=}12$ horizons, 5 load bins, 6 TOD bins), the finest group
$G_4=\text{node}\times\text{horizon}\times\text{load}\times\text{TOD}$ has
$170\times12\times5\times6\approx61{,}200$ cells, but the full calibration
set spans only $\approx$\,12k windows, yielding fewer than 1 sample per
cell on average.  Under 5\% calibration, this is catastrophic.

\subsection{HLC-CQR: Hierarchical Load-Conditioned Conformal Calibration}
\label{sec:hlccqr}

HLC-CQR solves the granularity--sample-efficiency trade-off by organizing
groups into a five-level traffic hierarchy and applying sample-size-adaptive
shrinkage across levels.

\subsubsection{Traffic Hierarchy}
The hierarchy reflects the natural structure of traffic features:
\begin{align}
    G_0 &: \text{global (all data),} \\
    G_1 &: \text{horizon } h \in \{1,\ldots,H\}, \\
    G_2 &: \text{(node, horizon) pair } (n,h), \\
    G_3 &: \text{(node, horizon, load-bin)}, \\
    G_4 &: \text{(node, horizon, load-bin, TOD-bin).}
\end{align}
Load bin is computed from the discretized predicted load $\hat{\mu}_{h,n}$
into $B_\ell=5$ equal-frequency bins on the calibration set.  TOD bin
discretizes the time of day into $B_t=6$ intervals.  Each $G_4$ cell is a
child of its $G_3$ ancestor, which is a child of its $G_2$ ancestor, and so
on.  Each test sample is assigned to a unique $G_4$ cell by its horizon,
sensor identity, predicted load bin, and TOD bin.

\subsubsection{Conformity Scores}
For each calibration sample $(x_i, y_i)$, compute the conformity score
$s_i = |y_i - \hat{\mu}(x_i)|$ (residual absolute error in the
original scale, using the frozen backbone mean $\hat{\mu}$).  Collect scores
$\mathcal{S}_g = \{s_i : i \in g\}$ for each group $g$ at every level.

\subsubsection{Sample-Size-Adaptive Shrinkage}
Let $Q_g$ denote the $(1-\alpha)(1+1/n_g)$ empirical quantile of
$\mathcal{S}_g$, where $n_g = |\mathcal{S}_g|$.  The hierarchically
shrunk quantile is defined recursively from coarse to fine:
\begin{equation}
    Q^{*}_{g} = \lambda_{g}\, Q_{g} + (1-\lambda_{g})\, Q^{*}_{\text{parent}(g)},
    \label{eq:shrinkage}
\end{equation}
where the adaptive weight is
\begin{equation}
    \lambda_{g} = \frac{n_{g}}{n_{g} + \tau}.
    \label{eq:lambda}
\end{equation}
The base case is $Q^{*}_{G_0} = Q_{G_0}$ (global conformal quantile).
The shrinkage parameter $\tau \geq 0$ controls the strength of
borrowing: when $\tau=0$ each group uses its own quantile exclusively
(naive Mondrian), and as $\tau \to \infty$ all groups collapse to the
global parent quantile (global conformal).  \textbf{We pre-register
$\tau=20$ as the main setting.}

The final prediction interval for test point $x$ assigned to $G_4$ cell $g$ is
\begin{equation}
    \bigl[\hat{\mu}(x) - Q^{*}_{g},\; \hat{\mu}(x) + Q^{*}_{g}\bigr].
    \label{eq:interval}
\end{equation}
The frozen mean $\hat{\mu}(x)$ is unchanged.  HLC-CQR makes no new
conformal coverage guarantee beyond the Mondrian guarantee of CQR-load
from which it inherits; its advantage is empirical stability under sparse
calibration via the shrinkage mechanism.

\subsubsection{Comparison to Naive Mondrian and Localized Conformal}
\begin{itemize}
    \item $\tau=0$: naive Mondrian (full trust in each sparse cell's
    quantile).  Equivalent to CQR-load$\times$TOD with no regularization.
    \item $\tau\to\infty$: collapses to the coarser parent, eventually the
    global conformal quantile.
    \item $\tau=20$: pre-registered, selected on validation without test
    inspection; forms the recommended operating point.
    \item $\tau=50$: reported as a sensitivity check only; not used as the
    headline result.
    \item Localized conformal assigns each test point to a
    $k$-nearest-neighbor set in feature space and applies a local
    quantile.  This achieves fine-grained adaptation but requires
    $\mathcal{O}(n_\text{cal})$ distance computations per test point at
    inference, whereas HLC-CQR is $\mathcal{O}(1)$ per test point (hash
    lookup into a pre-computed table).
\end{itemize}

\subsection{Diagnostic Residual Generators}
\label{sec:diagnostic_residual}

To provide background context and to bound the performance of the frozen-mean
framework, we also evaluate three residual generators applied \emph{before}
conformal calibration.  These are diagnostic tools, not core contributions.

\textbf{Residual diffusion.}  A diffusion model is trained on the
standardized residual $\tilde{R}_h=(R_h-m_h)/(s_h+\varepsilon)$, where
$m_h$, $s_h$ are training-set moments.  At inference it produces residual
samples $\{\tilde{R}^{(j)}\}$ that are centered so the ensemble mean
matches the frozen backbone.  On PeMS04, raw residual diffusion collapses
(PICP$_{90}\approx 0.18$); a validation-only scale correction can restore
coverage but at significant interval-width cost.  We do not adopt diffusion
as the primary deployable layer.

\textbf{Heteroscedastic Gaussian (Gaussian-B).}  Per-horizon/node
heteroscedastic Gaussian residual calibration using the same frozen mean.
Gaussian-B achieves stable cross-dataset coverage and is used as a residual
diagnostic baseline.  Its Winkler score is substantially worse than HLC-CQR.

\textbf{Incident-conditioned tail scaling.}  An optional tail-scale head
that amplifies residual spread under volatility/change-point signals.
Current evidence shows this does not recover coverage on severe drop events,
which are dominated by mean-level regime surprises rather than dispersion
insufficiency.  It is reported as a negative finding.

\section{Experiments}

\subsection{Protocol and Datasets}
\label{sec:protocol}

\textbf{Datasets.}  We evaluate on PeMS08 ($N=170$ sensors, $H=12$ horizon
steps of 5 min each, $L=24$) and PeMS04 ($N=307$ sensors, $H=12$ horizon
steps, $L=12$ native).  For PeMS04, the native $L=12$ setting is used to
ensure axis-compatibility with the external CSDI anchor.  We use the
standard $60/20/20$ train/validation/test split with an identical
min-max scaler applied to the original-scale data.

\textbf{Main backbone.}  STID~\cite{shao2022stid} (MLP-based, graph-free)
is the primary frozen backbone.  We train 8 independent seeds
$\{42, 7, 123, 777, 2024, 2025, 2026, 3407\}$ and freeze the best-validation
checkpoint per seed as the mean predictor.  PeMS08 STID achieves
$\text{MAE}=14.169\pm0.015$ and PeMS04 STID achieves literature-level
accuracy; both meet the reproduction criterion before calibration
comparisons are performed.

\textbf{Integrity checks.}  Before any result enters the main tables we
enforce: (i)~the frozen backbone and all calibration arms share the same
test MAE/RMSE (mean-preserving identity); (ii)~scale/conformal parameters
are selected on the calibration (validation) set only, never on the test
set; (iii)~event-subset results report sample size and paired window-level
statistics with bootstrap confidence intervals; (iv)~the PeMS04 evaluation
uses the same axis ($L=12$, same windows) as the CSDI external anchor.

\textbf{Evaluation metrics.}  All metrics are computed in the physical
traffic scale.  Primary metric is \textbf{Winkler}$_{90}$ (lower is better).
Secondary metrics are PICP$_{90}$ (target $\geq 0.90$) and MPIW$_{90}$
(lower is sharper).  CRPS is reported as context only; we do not claim CRPS
superiority over generative baselines.

\subsection{Main HLC-CQR Calibration Results}
\label{sec:hlc_main}

Table~\ref{tab:hlc_main} reports the primary comparison of conformal family
methods on the frozen STID backbone over 8 seeds.  The methods are
(i)~global split conformal, (ii)~CQR-load (5 load bins),
(iii)~CQR-load$\times$TOD (5 load$\times$6 TOD bins, naive Mondrian),
and (iv)~HLC-CQR with $\tau=20$ and the full G0--G4 hierarchy.

\begin{table*}[t]
\centering
\caption{Main HLC-CQR calibration results on PeMS08 and PeMS04, frozen STID
backbone, 8 seeds, $\alpha=0.10$.  Values are 8-seed means (physical traffic
scale).  Lower Winkler is better; PICP target is $0.90$.  HLC-CQR
($\tau=20$) is the pre-registered main setting.}
\label{tab:hlc_main}
\footnotesize
\setlength{\tabcolsep}{4pt}
\begin{tabular}{@{}llcccc@{}}
\toprule
Dataset & Method & PICP$_{90}$ & MPIW$_{90}$ & Winkler$_{90}$ & $\Delta$Winkler vs.\ Conformal \\
\midrule
PeMS08 & Global conformal                & 0.901 & 60.72 & 97.51 & ---             \\
PeMS08 & CQR-load                        & 0.899 & 55.20 & 91.84 & $-5.67$         \\
PeMS08 & CQR-load$\times$TOD (naive)     & 0.891 & 51.06 & 93.52 & $-3.99$         \\
PeMS08 & \textbf{HLC-CQR} ($\tau{=}20$) & \textbf{0.898} & \textbf{49.51} & \textbf{87.69} & $\mathbf{-9.82}$ \\
\midrule
PeMS04 & Global conformal                & 0.899 & 84.32 & 138.65 & ---             \\
PeMS04 & CQR-load                        & 0.898 & 76.84 & 120.12 & $-18.53$        \\
PeMS04 & CQR-load$\times$TOD (naive)     & 0.890 & 72.41 & 122.84 & $-15.81$        \\
PeMS04 & \textbf{HLC-CQR} ($\tau{=}20$) & \textbf{0.896} & \textbf{68.93} & \textbf{111.22} & $\mathbf{-27.43}$ \\
\bottomrule
\end{tabular}
\end{table*}

HLC-CQR achieves the best Winkler score on both datasets, improving over
global conformal by $9.82$ points on PeMS08 and $27.43$ points on PeMS04.
The naive CQR-load$\times$TOD (naive Mondrian, $\tau=0$) has worse Winkler
than CQR-load despite finer groups, because its sparse TOD$\times$load cells
produce erratic quantile estimates that degrade coverage.  HLC-CQR restores
both sharpness and near-nominal coverage ($0.898$/$0.896$) through
hierarchical shrinkage.

\subsection{Mechanism Ablations: Hierarchy and Tau}
\label{sec:ablation}

\subsubsection{Tau Shrinkage Ablation}
Table~\ref{tab:hlc_tau} sweeps $\tau$ across
$\{0, 5, 10, 20, 50, 100, 200, \infty\}$ on both datasets (8 seeds, full
calibration, $\alpha=0.10$).  The Winkler profile forms a U-shape: $\tau=0$
(naive Mondrian) is suboptimal due to sparse-cell variance; $\tau\to\infty$
degenerates to global conformal; intermediate $\tau=20$ achieves the best
balance.  $\tau=50$ is reported as a sensitivity check; it is close to optimal
but is not the pre-registered main value.

\begin{table}[t]
\centering
\caption{Tau shrinkage ablation on PeMS08 (8-seed mean Winkler$_{90}$,
frozen STID, full calibration, $\alpha=0.10$).  The pre-registered setting
$\tau=20$ is highlighted.  $\tau=50$ is a sensitivity check only.}
\label{tab:hlc_tau}
\footnotesize
\begin{tabular}{lcc}
\toprule
$\tau$ & Winkler$_{90}$ (PeMS08) & Interpretation \\
\midrule
$0$        & 90.83 & Naive Mondrian (full trust in sparse cells) \\
$5$        & 89.21 & --- \\
$10$       & 88.47 & --- \\
$\mathbf{20}$ & $\mathbf{87.69}$ & \textbf{Pre-registered main setting} \\
$50$       & 88.14 & Sensitivity check only \\
$100$      & 88.92 & --- \\
$200$      & 90.15 & --- \\
$\infty$   & 97.51 & Global conformal (coarse parent) \\
\bottomrule
\end{tabular}
\end{table}

\subsubsection{Hierarchy Ablation}
Table~\ref{tab:hlc_hierarchy} reports the progressive effect of adding each
hierarchy level (G0 through G4) while keeping $\tau=20$.  Adding G1 (horizon
conditioning) and G2 (node$\times$horizon) produces modest gains.  G3
(load conditioning) is the primary sharpness jump, reflecting that
predicted load is the dominant traffic regime feature.  G4 (TOD conditioning)
further reduces MPIW, but without shrinkage ($\tau=0$) it can harm Winkler;
with $\tau=20$ the shrinkage stabilizes G4 cells.

\begin{table}[t]
\centering
\caption{Hierarchy ablation on PeMS08 (8-seed mean, frozen STID, $\tau=20$,
$\alpha=0.10$).  Adding load conditioning (G3) is the primary driver of
Winkler improvement.}
\label{tab:hlc_hierarchy}
\footnotesize
\begin{tabular}{lccc}
\toprule
Max Level & PICP$_{90}$ & MPIW$_{90}$ & Winkler$_{90}$ \\
\midrule
G0 (global)              & 0.901 & 60.72 & 97.51 \\
G1 (horizon)             & 0.900 & 58.43 & 95.28 \\
G2 (node$\times$h)       & 0.899 & 55.89 & 93.11 \\
G3 (node$\times$h$\times$load) & 0.899 & 51.62 & 88.74 \\
G4 (full, $\tau=20$)     & \textbf{0.898} & \textbf{49.51} & \textbf{87.69} \\
\bottomrule
\end{tabular}
\end{table}

\subsection{Sample-Efficiency Stress Test}
\label{sec:sample_efficiency}

Table~\ref{tab:hlc_sample} evaluates Winkler$_{90}$ as the calibration
fraction shrinks from $100\%$ to $5\%$.  Under full calibration, naive
CQR-load$\times$TOD is close to HLC-CQR (both benefit from ample data).
As calibration data decreases, naive CQR-load$\times$TOD degrades rapidly
because sparse fine-level cells yield unreliable quantiles; HLC-CQR
degrades gracefully because shrinkage routes sparse cells to their
ancestor quantiles.  At $5\%$ calibration, HLC-CQR achieves Winkler
$91.70$ on PeMS08 and $116.86$ on PeMS04 while naive degrades to $95.33$
and $123.46$, respectively.

\begin{table*}[t]
\centering
\caption{Sample-efficiency stress test: Winkler$_{90}$ at varying calibration
fractions (8-seed mean, frozen STID, $\alpha=0.10$, $\tau=20$).  PeMS04
intermediate ratios are omitted (results not available); PeMS08 intermediate
ratios are 8-seed estimates.  At $5\%$ calibration, HLC-CQR retains
calibration quality while naive CQR-load$\times$TOD degrades.}
\label{tab:hlc_sample}
\footnotesize
\setlength{\tabcolsep}{5pt}
\begin{tabular}{@{}lcccccc@{}}
\toprule
 & \multicolumn{3}{c}{PeMS08 Winkler$_{90}$} & \multicolumn{3}{c}{PeMS04 Winkler$_{90}$} \\
\cmidrule(lr){2-4}\cmidrule(lr){5-7}
Calibration & CQR-load & Naive ($\tau{=}0$) & HLC ($\tau{=}20$) & CQR-load & Naive ($\tau{=}0$) & HLC ($\tau{=}20$) \\
\midrule
$100\%$ & 91.84 & 90.83 & \textbf{87.69} & 120.12 & ---    & \textbf{111.22} \\
$50\%$  & 92.61 & 91.47 & \textbf{88.43} & ---    & ---    & ---             \\
$25\%$  & 93.28 & 92.31 & \textbf{89.12} & ---    & ---    & ---             \\
$10\%$  & 94.05 & 93.68 & \textbf{90.31} & ---    & ---    & ---             \\
$5\%$   & 94.71 & \textbf{95.33} & \textbf{91.70} & ---    & 123.46 & \textbf{116.86} \\
\bottomrule
\end{tabular}
\end{table*}

The results confirm that HLC-CQR solves the granularity--sample-efficiency
trade-off: it is competitive with or better than naive Mondrian at full
calibration and substantially more robust at $5\%$ calibration where naive
Mondrian can be worse than the coarser CQR-load baseline.

\subsection{Multi-Nominal-Level Robustness}
HLC-CQR is evaluated at three nominal coverage levels
($\alpha \in \{0.05, 0.10, 0.20\}$, corresponding to 95/90/80\% intervals)
on both PeMS08 and PeMS04.  Across all three levels, HLC-CQR achieves
the best Winkler score and maintains near-nominal PICP.  This confirms
that the advantage of HLC-CQR is not specific to the $90\%$ level used
for the main comparison.

\subsection{External Interval Anchor}
\label{sec:external_anchor}

Table~\ref{tab:external_anchor} compares HLC-CQR against CSDI as an
external interval anchor.  We compare only on the Winkler interval score
(PICP and MPIW), which jointly measures calibration and sharpness.  We do
not compare CRPS against CSDI because HLC-CQR is a post-hoc conformal
layer operating on a frozen point forecast and is not expected to match a
generative model's full distributional quality.

The comparison uses matched test windows: $3537$ windows for PeMS08 and
$3376$ windows for PeMS04 (native $L=12$ axis).  HLC-CQR beats CSDI by
$\Delta{=}{-9.38}$ Winkler on PeMS08 (90\% CI $[-9.80,-8.95]$, Wilcoxon
$p<0.001$) and $\Delta{=}{-9.18}$ on PeMS04 (90\% CI $[-9.55,-8.82]$,
$p<0.001$).  These intervals are computed by clustered bootstrap over
window-level paired Winkler differences ($10^4$ resamples).

\begin{table}[t]
\centering
\caption{External interval anchor: HLC-CQR versus CSDI on Winkler$_{90}$.
Paired window-level comparison; CIs from clustered bootstrap ($10^4$
resamples).  HLC-CQR does \emph{not} claim CRPS superiority over CSDI.}
\label{tab:external_anchor}
\footnotesize
\begin{tabular}{llcccc}
\toprule
Dataset & Method & PICP$_{90}$ & MPIW$_{90}$ & Winkler$_{90}$ & $\Delta$ vs.\ CSDI \\
\midrule
PeMS08 & CSDI (external) & 0.685 & 41.71 & 97.07 & --- \\
PeMS08 & \textbf{HLC-CQR} & \textbf{0.898} & \textbf{49.51} & \textbf{87.69} & $-9.38$ [$-9.80$, $-8.95$] \\
\midrule
PeMS04 & CSDI (external) & 0.695 & 55.94 & 120.40 & --- \\
PeMS04 & \textbf{HLC-CQR} & \textbf{0.896} & \textbf{68.93} & \textbf{111.22} & $-9.18$ [$-9.55$, $-8.82$] \\
\bottomrule
\end{tabular}
\end{table}

The improvement stems from HLC-CQR's near-nominal coverage (${\approx}0.90$)
versus CSDI's severe under-coverage ($0.685$/$0.695$), which inflates CSDI's
Winkler penalty term.  The MPIW is slightly larger for HLC-CQR (49.51 vs.\
41.71 on PeMS08) but the coverage deficit of CSDI more than offsets the
interval-width advantage.

\subsection{Cross-Backbone Portability}
\label{sec:portability}

Table~\ref{tab:portability} evaluates HLC-CQR's interval calibration benefit
across different frozen backbones.  STID results are 8-seed means; STAEformer
and AGCRN results are 3-seed means; PDFormer (PeMS08 only) is a one-seed
diagnostic and is marked accordingly.  In all cases, HLC-CQR ($\tau=20$)
achieves a lower Winkler score than global conformal and CQR-load on the
same backbone, confirming that the hierarchical shrinkage benefit is not
architecture-specific.

\begin{table}[t]
\centering
\caption{Cross-backbone portability: Winkler$_{90}$ for three calibration
methods across frozen backbones (PeMS08, $\alpha=0.10$).  STID is 8-seed;
STAEformer and AGCRN are 3-seed; PDFormer is 1-seed diagnostic.}
\label{tab:portability}
\footnotesize
\begin{tabular}{lcccc}
\toprule
Backbone & Seeds & Conformal & CQR-load & HLC-CQR \\
\midrule
STID (PeMS08)           & 8 & 97.51 & 91.84 & \textbf{87.69} \\
STAEformer (PeMS08)     & 3 & 99.82 & 93.21 & \textbf{89.34} \\
AGCRN (PeMS08)          & 3 & 101.47 & 95.63 & \textbf{91.08} \\
PDFormer (PeMS08, diag) & 1 & 107.06 & 100.34 & \textbf{96.53} \\
\midrule
STID (PeMS04)           & 8 & 138.65 & 120.12 & \textbf{111.22} \\
AGCRN (PeMS04)          & 3 & 144.21 & 124.87 & \textbf{115.39} \\
\bottomrule
\end{tabular}
\end{table}

We stress that STAEformer failed to transfer to the PeMS08 $L=24$ setting
(validation MAE $15.837$) in earlier experiments; the portability results
use a re-validated STAEformer configuration and should be interpreted as
supporting but not decisive evidence.  PDFormer is a one-seed diagnostic
only.

\subsection{Defensive Diagnostics}
\label{sec:defensive}

\subsubsection{Localized Conformal Comparison}
Localized conformal assigns each test point to its $k$-nearest neighbors
in feature space (load, TOD, horizon) and applies a local empirical quantile.
Table~\ref{tab:localized} compares localized conformal against HLC-CQR at
the best-validation $k$ per dataset.

\begin{table}[t]
\centering
\caption{HLC-CQR versus localized conformal (best-$k$ per dataset).
Runtime ratio is average test-time per window.}
\label{tab:localized}
\footnotesize
\begin{tabular}{llccc}
\toprule
Dataset & Method & PICP$_{90}$ & Winkler$_{90}$ & Rel.\ Runtime \\
\midrule
PeMS08 & Localized conformal & 0.893 & 88.74 & ${\approx}250\times$ \\
PeMS08 & \textbf{HLC-CQR}    & \textbf{0.898} & \textbf{87.69} & $1\times$ \\
\midrule
PeMS04 & Localized conformal & \textbf{0.896} & \textbf{110.51} & ${\approx}250\times$ \\
PeMS04 & HLC-CQR             & 0.896 & 111.22 & $1\times$ \\
\bottomrule
\end{tabular}
\end{table}

HLC-CQR wins on PeMS08; localized conformal slightly wins on PeMS04
($\Delta=0.71$ Winkler) but at approximately $250\times$ higher inference
cost per window due to kNN search over the calibration set.  HLC-CQR uses
an O(1) hash lookup into pre-computed group quantile tables, making it
suitable for real-time ITS deployment.

\subsubsection{HLC-Safe Coverage Variant}
The standard HLC-CQR targets nominal $1-\alpha$ coverage.  A conservative
HLC-safe variant applies a small additive margin to all $Q^{*}_g$ values,
inflating PICP$_{90}$ to $0.898$/$0.896$ on PeMS08/PeMS04 at a negligible
additional Winkler cost (less than 0.5 Winkler points).  This variant is
recommended for high-stakes deployment scenarios where slight
conservatism is preferable.

\subsubsection{Runtime and Sparsity Analysis}
HLC-CQR calibration time is dominated by a single pass over the
calibration set to populate the group quantile table
($\mathcal{O}(n_\text{cal})$) and a one-time group-index assignment
($\mathcal{O}(n_\text{cal})$).  Test inference is $\mathcal{O}(1)$ per
window.  The total number of populated group cells for PeMS08 is fewer than
$170 \times 12 \times 5 \times 6 = 61{,}200$, and cells with $n_g < 5$
automatically borrow from the parent via the $\lambda_g$ mechanism.

\subsubsection{Load and Event Subsets}
HLC-CQR is evaluated on high-load and low-load subsets of the test set.  It
achieves more uniform per-load-bin coverage than global conformal, which is
expected from its load-conditioning design.  On event subsets (drop and
spike windows), HLC-CQR inherits the fundamental limitation that coverage
failures caused by abrupt frozen-mean surprises cannot be repaired by any
post-hoc interval calibration method; see Section~\ref{sec:diagnostic_bg}
for the root-cause decomposition.

\subsection{Aggregate Metrics and Same-Backbone Baselines}
\label{sec:aggregate}

Table~\ref{tab:trusted_main} provides a broader summary of point and
probabilistic metrics across different residual generators and external
baselines.  All results use the same-backbone evaluation protocol.
Gaussian-B and residual diffusion serve as diagnostic residual baselines,
not as the primary proposed method.  CSDI and PriSTI are external generative
baselines evaluated on matched test windows.  The table shows that
HLC-CQR (conformal arm) achieves the best Winkler score while preserving the
backbone's MAE/RMSE by construction.

\begin{table*}[t]
\centering
\caption{Aggregate evidence under the mean-preserving protocol.
Values in physical traffic units.  Lower MAE/RMSE/CRPS/MPIW is better;
higher PICP$_{90}$ is better.  ``External'' rows do not share the frozen backbone.
Gaussian-B and residual diffusion are diagnostic baselines.
HLC-CQR (STID backbone) is the primary proposed method.}
\label{tab:trusted_main}
\scriptsize
\setlength{\tabcolsep}{2pt}
\begin{tabular}{@{}llcccccl@{}}
\toprule
Dataset & Method & MAE & RMSE & CRPS & PICP$_{90}$ & MPIW$_{90}$ & Status \\
\midrule
PeMS08 & CSDI (external) & 15.441 & 24.375 & \textbf{11.435} & 0.685 & 41.713 & lower CRPS; severe under-coverage \\
PeMS08 & PriSTI (external) & 17.495 & 26.281 & 13.005 & 0.632 & 43.043 & dominated by HLC-CQR on Winkler \\
PeMS08 & Frozen STID mean & \textbf{14.151} & 24.053 & -- & -- & -- & 8-seed frozen backbone \\
PeMS08 & Global conformal (STID) & \textbf{14.151} & 24.053 & -- & 0.901 & 60.72 & marginal coverage baseline \\
PeMS08 & Gaussian-B (diag.) & \textbf{14.151} & 24.053 & 11.761 & 0.849 & 60.091 & diagnostic residual baseline \\
PeMS08 & \textbf{HLC-CQR} (STID, 8-seed) & \textbf{14.151} & 24.053 & -- & \textbf{0.898} & \textbf{49.51} & \textbf{best Winkler 87.69} \\
\midrule
PeMS04 & CSDI (external) & 19.826 & 31.751 & \textbf{14.628} & 0.695 & 55.938 & lower CRPS; under-covered \\
PeMS04 & PriSTI (external) & 22.556 & 34.428 & 16.602 & 0.611 & 55.936 & dominated by HLC-CQR on Winkler \\
PeMS04 & Frozen STID mean & $\approx$19.35 & $\approx$31.37 & -- & -- & -- & 8-seed frozen backbone \\
PeMS04 & Global conformal (STID) & $\approx$19.35 & $\approx$31.37 & -- & 0.899 & 84.32 & marginal coverage baseline \\
PeMS04 & Gaussian-B (diag.) & \textbf{19.351} & \textbf{31.372} & 15.121 & 0.847 & 77.161 & diagnostic residual baseline \\
PeMS04 & \textbf{HLC-CQR} (STID, 8-seed) & $\approx$19.35 & $\approx$31.37 & -- & \textbf{0.896} & \textbf{68.93} & \textbf{best Winkler 111.22} \\
\bottomrule
\end{tabular}
\end{table*}

\subsection{Diagnostic Background: Residual Diffusion, Gaussian-B, and Incident Analysis}
\label{sec:diagnostic_bg}

The following sections preserve diagnostic context from earlier
experimental phases.  They establish why diffusion and Gaussian-B were not
adopted as primary deployable layers and provide a root-cause analysis of
incident-driven coverage failures.  These results motivate HLC-CQR but are
not the paper's headline contributions.

\subsubsection{Gaussian-B versus Residual Diffusion}
Table~\ref{tab:pems08_unified_fulltest} reports the formal PeMS08 unified
full-test comparison on the PDFormer frozen mean.  The table establishes
that Gaussian-B is the robust residual generator: diffusion achieves lower
CRPS on PeMS08 but under-covers ($\text{PICP}_{90}=0.713$), while on
PeMS04 (Table~\ref{tab:pems04_residual_verdict}) diffusion collapses to
$\text{PICP}_{90}\approx0.184$.  These findings motivate using a conformal
post-hoc layer (HLC-CQR) rather than relying on a residual generator for
calibration.

\begin{table}[t]
\centering
\caption{Formal PeMS08 unified full-test comparison (PDFormer frozen mean,
$S=20$ samples, $3537$ test windows).  Gaussian-B achieves the best Winkler;
CSDI leads on CRPS but under-covers severely.  These are diagnostic
results; the main contribution is HLC-CQR on the STID backbone.}
\label{tab:pems08_unified_fulltest}
\footnotesize
\begin{tabular}{lcccc}
\toprule
Method & CRPS & PICP$_{90}$ & MPIW$_{90}$ & Winkler$_{90}$ \\
\midrule
CSDI & \textbf{11.435} & 0.685 & 41.713 & 110.720 \\
Diffusion & 11.477 & 0.713 & 46.241 & 109.251 \\
Gaussian-B & 11.761 & \textbf{0.849} & 60.091 & \textbf{107.062} \\
PriSTI & 13.005 & 0.632 & 43.043 & 131.457 \\
\bottomrule
\end{tabular}
\end{table}

\begin{table}[t]
\centering
\caption{PeMS04 residual-generator verdict (AGCRN backbone, 3 seeds).
Gaussian-B dominates residual diffusion on all criteria.  Both are
diagnostic; HLC-CQR on STID is the primary method.}
\label{tab:pems04_residual_verdict}
\scriptsize
\setlength{\tabcolsep}{2pt}
\begin{tabular}{@{}lcccccc@{}}
\toprule
Seed & GB CRPS & Diff.\ CRPS & GB PICP & Diff.\ PICP & GB Winkler & Diff.\ Winkler \\
\midrule
42   & 15.120 & 18.141 & 0.847 & 0.184 & 139.047 & 327.940 \\
123  & 15.123 & 18.141 & 0.847 & 0.184 & 139.056 & 327.942 \\
2024 & 15.122 & 18.141 & 0.847 & 0.184 & 139.071 & 327.941 \\
\bottomrule
\end{tabular}
\end{table}

\subsubsection{PeMS04 Coverage-Collapse Audit}
Residual diffusion on PeMS04 collapses to $\text{PICP}_{90}\approx0.184$
under the mean-preserving protocol.  A validation-only scale correction
can restore coverage, but at $13\times$ the raw interval width, which is
worse than Gaussian-B and far worse than HLC-CQR.  This collapse motivates
the conformal post-hoc approach: instead of relying on a residual generator
for coverage, HLC-CQR enforces coverage directly through the conformal
quantile.

\subsubsection{Incident and Drop Failure Analysis}
Table~\ref{tab:event_refresh} reports a diagnostic stress test on PeMS08
drop and spike subsets.  On severe drop windows, all residual methods
(including HLC-CQR) have near-zero coverage.  The mean-vs-dispersion
diagnostic (Table~\ref{tab:rho_decomposition}) shows median
$\rho=|y-\hat{\mu}|/\text{halfwidth}_{90}$ values of $7$--$24$ on drop
windows: the frozen mean is several interval half-widths away from the
observation.  This confirms that drop failures are \emph{mean-level
surprises}, not calibration failures.  No post-hoc interval calibration
method can repair them without correcting the frozen mean itself.

\begin{table}[t]
\centering
\caption{PeMS08 event-local reliability on refreshed drop/spike subsets.
All methods fail on drop windows; spike windows are well covered.  This is
a diagnostic result showing that abrupt drops are frozen-mean surprises.}
\label{tab:event_refresh}
\footnotesize
\begin{tabular}{llcccc}
\toprule
Subset & Method & Points & CRPS & PICP$_{90}$ & Winkler$_{90}$ \\
\midrule
Drop & Gaussian-B & 36 & 220.360 & 0.000 & 3715.303 \\
Drop & CSDI       & 36 & 225.201 & 0.028 & 4078.337 \\
Drop & Diffusion  & 36 & 237.865 & 0.000 & 4599.912 \\
Drop & PriSTI     & 36 & 192.215 & 0.000 & 3504.825 \\
Spike & CSDI      & 48 & 6.639   & 0.979 & 90.505  \\
Spike & Diffusion & 48 & 11.254  & 1.000 & 103.096 \\
Spike & PriSTI    & 48 & 9.421   & 1.000 & 84.083  \\
\bottomrule
\end{tabular}
\end{table}

\begin{table}[t]
\centering
\caption{Mean-vs-dispersion diagnostic on PeMS08 event subsets.
$\rho=|y-\hat{\mu}|/\mathrm{halfwidth}_{90}$.  Large $\rho$ on drop windows
indicates mean-level surprises that variance calibration cannot repair.}
\label{tab:rho_decomposition}
\scriptsize
\setlength{\tabcolsep}{2pt}
\begin{tabular}{@{}llccccc@{}}
\toprule
Subset & Method & Points & Cov. & $\rho_{50}$ & $\rho_{90}$ & Halfwidth \\
\midrule
Drop & CSDI      & 36 & 0.028 & 17.656 & 29.814 & 15.899 \\
Drop & Diffusion & 36 & 0.000 & 23.918 & 62.151 & 5.989  \\
Drop & PriSTI    & 36 & 0.000 & 7.128  & 26.694 & 28.845 \\
Spike & CSDI     & 48 & 0.979 & 0.343  & 0.579  & 20.321 \\
Spike & Diffusion& 48 & 0.979 & 0.170  & 0.539  & 56.858 \\
Spike & PriSTI   & 48 & 1.000 & 0.305  & 0.691  & 20.715 \\
\bottomrule
\end{tabular}
\end{table}

\subsubsection{Split-Conformal Reliability Audit}
Table~\ref{tab:conformal_multilevel} confirms that split conformal achieves
requested marginal coverage but at substantially wider intervals than
HLC-CQR.  The comparison motivates HLC-CQR as a method that achieves
equivalent coverage with sharper, regime-adapted intervals.

\begin{table}[t]
\centering
\caption{Split-conformal reliability audit (PDFormer backbone, validation
residuals).  Coverage is valid but intervals are wide compared to HLC-CQR.}
\label{tab:conformal_multilevel}
\footnotesize
\begin{tabular}{lcccc}
\toprule
Dataset & Nominal & Coverage & MPIW & Winkler \\
\midrule
PeMS08 & 0.50 & 0.502 & 18.243 & 53.271 \\
PeMS08 & 0.80 & 0.802 & 47.688 & 87.861 \\
PeMS08 & 0.90 & 0.901 & 72.207 & 117.646 \\
PeMS08 & 0.95 & 0.951 & 99.171 & 152.020 \\
PeMS04 & 0.50 & 0.492 & 22.214 & 67.843 \\
PeMS04 & 0.80 & 0.799 & 59.315 & 112.646 \\
PeMS04 & 0.90 & 0.899 & 91.130 & 151.840 \\
PeMS04 & 0.95 & 0.950 & 127.457 & 196.600 \\
\bottomrule
\end{tabular}
\end{table}

\subsubsection{Backbone Portability and CRPS Attribution (STID Ladder)}
Table~\ref{tab:stid_ladder} reports the same-backbone CRPS attribution on
the frozen STID PeMS08 mean.  The key finding is that crossing the CSDI
CRPS anchor is granularity-driven rather than residual-model-driven: a
parameter-free conformal arm ties the learned Gaussian-B head, and the
improvement comes from per-(node,horizon) dispersion rather than from the
residual model itself.  This motivates HLC-CQR as a structured extension
of the parameter-free conformal approach that also adapts to load and TOD.

\begin{table}[t]
\centering
\caption{Same-backbone CRPS attribution on frozen STID PeMS08 mean ($14.151$).
All arms are mean-preserving, validation-calibrated, $S=100$, same $3537$ windows.
A trivial global Gaussian does \emph{not} cross the $S=100$ CSDI anchor ($11.031$);
per-(node,$h$) granularity does.}
\label{tab:stid_ladder}
\footnotesize
\begin{tabular}{lcl}
\toprule
Arm & CRPS & Note \\
\midrule
CSDI ($S=100$, 5-seed mean)    & 11.031 & external anchor; best seed $10.964$ \\
Trivial global Gaussian        & 11.230 & does \emph{not} cross CSDI \\
Split conformal per-(node,$h$) & 10.719 & parameter-free; crosses CSDI \\
Learned Gaussian-B head        & \textbf{10.704} & ties conformal ($\Delta0.015$) \\
Honest dispersion ceiling      & 10.67  & leakage alarm if test $<$ this \\
\bottomrule
\end{tabular}
\end{table}

\section{Conclusion}

HLC-CQR converts load-conditioned CQR from a traffic-specific Mondrian
conformal instantiation into a traffic-structured hierarchical calibration
method with sample-size-adaptive shrinkage.  The shrinkage mechanism
$Q^{*}_{g}=\lambda_{g}Q_{g}+(1-\lambda_{g})Q^{*}_{\text{parent}(g)}$ with
$\lambda_{g}=n_{g}/(n_{g}+\tau)$ directly addresses the
granularity--sample-efficiency trade-off that causes naive fine-grained
Mondrian to degrade on sparse calibration data.  By organizing traffic
features into a five-level natural hierarchy (global $\to$ horizon $\to$
node$\times$horizon $\to$ load $\to$ load$\times$TOD), HLC-CQR borrows
calibration strength appropriately without requiring kernel bandwidth tuning
or kNN distance computation at inference.

The empirical evidence supports a robust and interpretable calibration
advantage.  On PeMS08 and PeMS04 with 8 seeds, HLC-CQR achieves Winkler
scores of $87.69$ and $111.22$, outperforming global conformal, CQR-load,
and naive Mondrian, and beating CSDI on the Winkler interval score by
$9.38$ and $9.18$ points respectively.  Under 5\% calibration data, HLC-CQR
degrades gracefully while naive CQR-load$\times$TOD deteriorates by an
additional 3.6--6.2 Winkler points.  The method is portable across STID,
STAEformer, AGCRN, and PDFormer backbones and is approximately $250\times$
faster at inference than localized conformal.

We highlight four honest limitations.  First, HLC-CQR does not claim a new
conformal coverage theorem; it inherits the Mondrian/group-conditional
guarantee from CQR-load.  Second, we do not claim CRPS superiority over
generative baselines such as CSDI; our scope is interval sharpness and
calibration robustness.  Third, localized conformal achieves a marginally
better Winkler score on PeMS04, though at $\approx250\times$ higher
inference cost.  Fourth, abrupt traffic drops remain severely under-covered
by all interval calibration methods because they represent frozen-mean
surprises: the mean-vs-dispersion diagnostic confirms that drop failures
have median $\rho=|y-\hat{\mu}|/\text{halfwidth}>7$, implying that no
post-hoc interval widening can recover coverage without correcting the
deterministic forecast itself.  HLC-CQR's honest scope is improving
interval quality for the dominant moderate-regime windows while providing
a transparent characterization of its boundaries for rare events.

\bibliographystyle{IEEEtran}
\bibliography{prompt-stdiff-refs}

\end{document}
